\settrack{trackbridge}
% VOICE: expository/pedagogical — teaches concepts, no narrative events, no "the COWS did X"

\chapter{The Factoring Game}
\label{spine:factoring-game}

% DMS MVP import from staging/raw/ — not final prose

\begin{quote}\small
\textit{``Privacy is necessary for an open society in the electronic age.''}

\hfill --- Eric Hughes, ``A Cypherpunk's Manifesto'' (1993)
\end{quote}

\vspace{0.5cm}

\srcnote{Cryptographic Origins.txt | staging/raw/pos09-the-factoring-game.md | 2026-02-15 | Full Cryptographic Origins document --- PKC explanation, crypto as munitions, DARPA need}

\section*{The State of Cryptography in 1990}
\label{spine:fg-the-state-of-cryptography-in-1990}

\srcnote{CH3 - Practical Cryptography \& project ULTRA II.txt | staging/raw/pos09-the-factoring-game.md | 2026-02-15 | State of Cryptography in 1990, ULTRA II project specs}

In 1990 the world's top spy organizations had an unhappy blind spot.  The mathematics of Public Key Cryptography, published publicly between 1976 and 1978, allowed civilian computer users to routinely encrypt communications in unbreakable ways.  Supposedly, officially, they still have this same blind spot in 2026.  Public key cryptography provides the S in HTTPS.  Alice and Bob could communicate privately.  Eve could analyse traffic, but she could not eavesdrop.  Organizations that existed to spy on communications, such as the United States National Security Agency and the other members of the Five Eyes, felt threatened.  The spy organizations badly needed a reliable way to crack public key cryptography. Something had to be done.

One need not be a mathematician to understand how Public Key Cryptography works.  It is based on the principle that it is harder to factor a number than to multiply its factors.  Quick, what are the prime factors of 91?  Stop reading while you calculate the answer.  After either answering that problem where you factor a number into its prime factors, or giving up, multiply thirteen by seven.  Which was harder?  Why?  The security of Public Key Cryptography comes from the mathematical difficulty of finding prime factors for very large numbers.  If the NSA spies wanted their cryptography problem solved, then this mathematical problem had to be solved.  This has been common knowledge to cryptographers since 1978.

In terms of US Government agencies in 1990, the obvious agency to try to solve the NSA's cryptography problem was DARPA.  DARPA is the Defense Advanced Research Projects Agency.  DARPA develops new technology for the military.  DARPA focuses on short-term (two to four year) projects run by small, purpose-built teams.  DARPA invented the internet circa 1969.

In order to crack public key cryptography DARPA needed to find a practical mathematical method for rapidly factoring large numbers.  This challenging topic has been the subject of many dissertations in mathematics.  A mathematician would call the desired process factoring in polynomial time.  While there are many tricks to speed the process, all of them are subject to exponential time slowdown as the size of the number increases.  Basically, all known approaches amount to a variation of trying every possible permutation.  Practical public/private key cryptography typically chooses a large enough number that it would take the best supercomputers many years of trial and error to find the prime factors and crack the key.\footcite{ferguson2010}

In 1992 David Deutsch and Richard Jozsa published the Deutsch-Jozsa algorithm.\footcite{deutschjozsa1992}  Although of little practical use, it is one of the first examples of a quantum algorithm that is exponentially faster than any possible deterministic classical algorithm.  Simon's problem (1994) demonstrated an exponential quantum speedup for a structured problem, inspiring Shor to apply the same periodicity-finding insight to integer factoring.  A quantum algorithm (software) capable of cracking public key cryptography was published by Peter Shor in 1994 --- \hovertiphtml{Shor's algorithm}.\footcite{shor1994}   The problem then, as now, is lack of appropriate quantum computation hardware on which to run such an algorithm.

By 1991 it was known to top cryptographers that quantum algorithms could, in theory, crack public key cryptography.  This was the basis for the 1992 Robert Redford movie \textit{Sneakers}, although quantum computation was not specifically mentioned.  No quantum hardware existed to run this quantum software.  The project planners at DARPA saw a possible way to solve their cryptanalysis problem.  They must invent a working quantum computation device.

\subsection*{The ULTRA II Specifications}
\label{spine:fg-the-ultra-ii-specifications}

The overall project specifications for ULTRA II were now clear:

Create a tool capable of efficiently cracking public key cryptography.   System must be able to rapidly determine the prime factors of a 4096 bit asymmetric key.  System must be extensible to potentially handle much larger keys.  System must have an input mechanism to accept the number to be factored.  System must have an output mechanism to return the prime factors of the input.

Basic Method: Build some kind of quantum computer running some quantum algorithm.

Problem: No one knows how to build a device capable of meeting the specifications.  Several approaches are known that can produce `toy sized' quantum computers, but none can scale to the thousands of qubits required to factor real cryptographic keys.  The main problem is \hovertiphtml{decoherence} time.

Everything above is public-domain history. The cryptographic arms race, DARPA's mandate, Shor's algorithm, and the decoherence problem are documented in the academic and popular literature. What follows is where the book's specific claims begin.

\vspace{1cm}\noindent\rule{0.5\textwidth}{0.5pt}\vspace{0.5cm}

This would not be the first time a computational breakthrough was built in secret, used to decisive effect, and then buried. At Bletchley Park in 1943--44, \deeplink{colossus-precedent}Tommy Flowers built Colossus --- a programmable electronic computer, operational years before ENIAC, designed to break the Lorenz cipher that carried Hitler's strategic communications. Colossus worked. Ten were built. After the war, on government orders, eight were destroyed. The remaining two were classified and dismantled in the late 1950s and 1960s. The official history of computing began with ENIAC in 1945. The actual history began at least two years earlier, in a room that officially didn't exist, solving a problem that officially couldn't be discussed. Colossus was not declassified until the mid-1970s. The full technical details were not released until 2000. For more than half a century, the people who built and operated the first real computer could not say so. Some died without ever being acknowledged. The GCHQ precedent that follows is another instance of the same pattern. It is worth understanding that the pattern is old.

\section*{The GCHQ Precedent}
\label{spine:fg-the-gchq-precedent}

\srcnote{HEALER-RECONSTRUCTION.md | staging/raw/pos09-the-factoring-game.md | 2026-02-15 | GCHQ/Cocks precedent and Crypto Wars --- classified independent discovery of PKC}

Shor's algorithm follows the GCHQ trajectory. First presented in April 1994 at Bell Labs, with a public talk at Cornell in May, formally published in November 1994 at FOCS. Shor was at AT\&T Bell Labs --- AT\&T had deep NSA partnerships, including BLARNEY (1978) and FAIRVIEW (1985). The NSA contacted Shor after his first talk and awarded him the Mathematics in Cryptology Award in 1995 --- rapid recognition for a supposedly ``surprise'' result.

In November 1996, \deeplink{shor-and-nsa}Executive Order 13026 moved encryption from the Munitions List to the Commerce Control List.\footcite{eo13026} The NSA had fought PKC export for twenty years (1976--1996), then reversed remarkably quickly --- just two years after Shor's algorithm. Under the narrative presented in this book, the NSA already had quantum cryptanalysis capability by 1995. Shor's public algorithm confirmed the direction. EO 13026 released the export restrictions because PKC was no longer worth protecting.

The proposition: the DARPA team discovered a variant of quantum factoring \textit{before} Shor's 1994 publication. This is consistent with the GCHQ/Cocks precedent --- independent classified discovery years before public discovery. The precise mechanism is unknown.

The GCHQ/Cocks precedent is real: classified independent discovery, years before public discovery, is documented history. The question is whether this pattern extends to the specific claims of this book. Under Possibility~A, the precedent is being used to make an unfalsifiable argument --- ``they did it before, so they could have done it again'' is true but proves nothing about whether they actually did. Under Possibility~B, DARPA may have had a quantum computing program in the early 1990s (it did), and some variant of quantum factoring may have been explored, but the leap to a fully operational system goes beyond what the GCHQ analogy supports. Under Possibility~C, the pattern is exact: classified discovery preceded public discovery by years, and EO~13026 is the evidence that something changed.

\chapterreturn
